\scanpage{22}

$\unsure{M_\bullet}\sim\unsure{\widetilde M_\bullet}$ en sens inverse \inserted{(ils sont adjoints...)} et le composé \deleted{est} $E\sim\widetilde{\underline{\Gamma}(E)}$ est isom. à l'identité, ce qui montre que $\underline{\Gamma}^\bullet$ est \uline{pleinement fidèle}.

$F$ correspond à un sous-\deleted{$\unsure{\mathcal O}$} $\mathcal A^\bullet$-module $\Gamma^\bullet(F)$ de $\Gamma^\bullet(E)$, pas à priori stable par action de $G$. Si \deleted{$F'\supset F$ est} \inserted{$F'$, tel que} $F\subset F'\subset E$, \deleted{$F'\supset F$}, est stable par $G$, alors $\Gamma^\bullet(F')$ \deleted{$\Gamma^\bullet_{\deleted{*}}(\unsure{E'|F'})\subset\Gamma^\bullet(E)$}, satisfait
\[
\underline{\Gamma}^\bullet(F)\subset\underline{\Gamma}^\bullet(F')\subset\underline{\Gamma}^\bullet(E)\quad\text{est est stable par }G.
\]

On sait qu'il y a un plus petit sous-\inserted{$\unsure{\underline O_S|\mathcal O_S}$}-Module $N$ de $\Gamma^\bullet(E)$, \unsure{contt} $\underline{\Gamma}^\bullet(F)$, et stable par $G$. On voit tout de suite que $N$ est un sous-module gradué [\unsure{par ex.} $N\cap\Gamma^i(E)$ contient $\underline{\Gamma}^i(F)$, et est stable par $G$, donc contt le plus petit sous-module $M^i$ de $\underline{\Gamma}^i(E)$ contt $\underline{\Gamma}^i(F)$ et stable par $G$, donc $N$ contient $\sum M^i$, d'où par minimalité lui est égal]. De $=$, par fonctorialité de la construction de $N$, on voit que $N$ est un sous-$\mathcal A$-module. Considérons alors

\mspage{21}
